\begin{pastexam}[label={px:ncu-113-dsa-16}]{中央資工所 113 學年度 第 16 題}
（複選；每題 5 分，全部答對才給分，答錯倒扣一分）

Which of the following statements about sorting algorithms are incorrect?
\begin{enumerate}[label=(\Alph*)]
\item The merge sort and quick sort are designed based on the similar algorithm design principle.
\item If most of the numbers are in the correct order, quick sort is the best algorithm.
\item The worst-case time complexity of quick sort, merge sort, and heap sort is $O(n\log n)$.
\item The worst-case time complexity and the expected time complexity of both merge sort and heap sort are in
the same order.
\item No sorting algorithm has time complexity lower than $O(n\log n)$.
\end{enumerate}
\end{pastexam}

\begin{pastsoln}
題目問的是\textbf{錯誤}的敘述。 答案：(B)(C)(E)。
\begin{itemize}
\item (A) 正確，不選。 合併排序和快速排序都是分治法：切成小問題、遞迴解、再組合起來。 差別只在主要的
工作放在哪裡：合併排序在遞迴之後合併，快速排序在遞迴之前分割（資結書第~14.5~節開頭）。
\item (B) 錯誤，要選。 數字大部分已經在正確位置時，逆序對個數 $I$ 很小，插入排序的時間是
$\Theta(n + I)$，接近線性（資結書第~14.1~節）。 快速排序若固定取第一個或最後一個當基準，排好的輸入正是
它的最壞情況：資結書第~14.5~節實際數過，$2000$ 個排好的數比了 $1{,}999{,}000$ 次。 改成隨機基準也還是
$\Theta(n\lg n)$，比不上插入排序。
\item (C) 錯誤，要選。 合併排序、堆積排序最壞是 $O(n\lg n)$，快速排序最壞是 $\Theta(n^2)$（每次都切成
$0$ 和 $m - 1$ 個）。
\item (D) 正確，不選。 合併排序不管輸入如何都是 $\Theta(n\lg n)$；堆積排序最壞 $O(n\lg n)$，鍵互不相同時
平均也是 $\Theta(n\lg n)$。 兩者的最壞和平均同一級（資結書第~14.7~節的比較表）。
\item (E) 錯誤，要選。 這句用 $O$ 寫，意思是「沒有任何排序比 $n\log n$ 快」。 $\Omega(n\lg n)$ 的下界只管
比較排序（資結書第~15.1~節）；鍵是 $0$ 到 $k$ 的整數、$k = O(n)$ 時，計數排序是 $\Theta(n)$。 就算只看
比較排序，插入排序在排好的輸入上也只要 $\Theta(n)$。
\end{itemize}
\end{pastsoln}
